\documentclass[11pt,a4paper,twoside]{article}

% =========================
% Packages
% =========================
\usepackage[margin=1in]{geometry}
\usepackage[T1]{fontenc}
\usepackage[utf8]{inputenc}
\usepackage{lmodern}
\usepackage{microtype}
\usepackage{setspace}
\usepackage{amsmath,amssymb,amsfonts,mathtools}
\usepackage{graphicx}
\usepackage{float}
\usepackage{booktabs}
\usepackage{tabularx}
\usepackage{array}
\usepackage{multirow}
\usepackage{adjustbox}
\usepackage{xcolor}
\usepackage{hyperref}
\usepackage{caption}
\usepackage{fancyhdr}

% =========================
% Formatting
% =========================
\setstretch{1.1}
\captionsetup{font=small,labelfont=bf}
\hypersetup{colorlinks=true, linkcolor=black, citecolor=black, urlcolor=blue}
\setlength{\parindent}{1.2em}
\setlength{\parskip}{0.4em}
\raggedbottom
\graphicspath{{../../reports/figures/paper/}{./}}

% =========================
% Header / Footer
% =========================
\pagestyle{fancy}
\fancyhf{}
\fancyhead[RO]{\small Solar PV Fault Diagnosis\hspace{1em}\thepage}
\fancyhead[LE]{\small\thepage\hspace{1em}Nguyen et al.}
\fancyfoot{}
\renewcommand{\headrulewidth}{0pt}
\renewcommand{\footrulewidth}{0pt}

\fancypagestyle{plain}{
    \fancyhf{}
    \fancyhead[RO]{\small Solar PV Fault Diagnosis\hspace{1em}\thepage}
    \fancyhead[LE]{\small\thepage\hspace{1em}Nguyen et al.}
    \fancyfoot{}
    \renewcommand{\headrulewidth}{0pt}
    \renewcommand{\footrulewidth}{0pt}
}

% =========================
% Title
% =========================
\title{\textbf{An Explainable, Data-Efficient Framework for Multi-Class Fault Diagnosis and Carbon Loss Quantification in Solar PV Monitoring}}

\author{%
\begin{tabular}{c}
Linh Hoang Nguyen\textsuperscript{1},
Tien Dong Dinh\textsuperscript{1,*},
An Dang Thai Le\textsuperscript{1},
Phuc Nguyen Huu Do\textsuperscript{1} \\[0.35em]
\textsuperscript{1}FPT University, Ho Chi Minh Campus, Vietnam \\[0.25em]
\small \texttt{LinhNH67@fe.edu.vn}; \texttt{tienddse200934@fpt.edu.vn}; \texttt{anldtse200518@fpt.edu.vn}; \texttt{do.huuphuc2k5@gmail.com} \\[0.20em]
\end{tabular}%
}
\date{}

\begin{document}

\maketitle
\thispagestyle{empty}

% =========================
% Abstract
% =========================
\begin{abstract}
Solar photovoltaic (PV) monitoring faces a fundamental tension: the per-inverter local models that achieve the highest forecasting accuracy inadvertently mask the very faults they should reveal, learning the degraded state of an underperforming inverter as its new normal. Moreover, existing approaches often frame fault diagnosis as a binary problem, ignoring the specific operational and maintenance (O\&M) actions required. This paper proposes an explainable, data-efficient framework for multi-class fault diagnosis and carbon loss quantification that inverts the conventional design objective: rather than fitting each inverter as closely as possible, we train a single global Expected Power Model, using Extreme Gradient Boosting, whose deviations expose degraded assets against an ideal plant-wide baseline. Combined with a physics-informed rule-based taxonomy, this design overcomes the anomaly-masking problem inherent to local architectures. We further introduce a Masked Loss Aggregation technique to eliminate Positive Drift Bias in energy loss quantification, enabling faults to be mapped to carbon-emission losses. Evaluated through physics-informed synthetic fault injection on a short (34-day) public benchmark, the framework achieves an F1 score of 0.74 $\pm$ 0.05, and identifies over 639 MWh of recoverable energy loss corresponding to nearly 450 metric tons of avoided CO\textsubscript{2}. The workflow demonstrates that practical inverter-level performance monitoring can be established from short-term operational data alone, without the multi-month or multi-year records required by many long-term degradation studies.
\end{abstract}

\noindent\textbf{Keywords:} Solar Photovoltaic, Fault Diagnosis, XGBoost, Anomaly Detection, Carbon Footprint, SHAP.

\vspace{1em}
\noindent\textbf{Abbreviations:} PV (Photovoltaic), O\&M (Operations and Maintenance), XGBoost (Extreme Gradient Boosting), SHAP (SHapley Additive exPlanations), UML (Unsupervised Machine Learning), LOF (Local Outlier Factor), IF (Isolation Forest), SVM (Support Vector Machine), CEA (Central Electricity Authority).
\vspace{1em}

% =========================
% 1. Introduction
% =========================
\section{Introduction}

Machine learning has been widely applied to PV power forecasting, where the objective is to predict generation as accurately as possible. Achieving this typically favors per-inverter local models that fit each unit's historical behavior closely; for instance, Nahar et al.~\cite{nahar2025} report high forecasting accuracy by training an individual model for every inverter. However, an objective that is ideal for forecasting is counterproductive for fault diagnosis. A model that fits a chronically underperforming inverter closely will learn its degraded output as the expected norm, producing near-zero residuals despite real power loss. We formally define this phenomenon as \textit{anomaly masking}: when an underperforming asset's historical data $X_{\text{local}}$ trains a model $f(X_{\text{local}})$ that normalizes the fault signature, yielding $y - f(X_{\text{local}}) \approx 0$ despite actual degradation. Traditional anomaly detection compounds this by framing diagnosis as a binary normal-versus-anomalous decision, providing limited utility to operations and maintenance (O\&M) teams who require actionable, multi-class categorizations to dispatch crews effectively.

This paper is built on a single observation: the design criterion for an expected-power model should be inverted when the goal shifts from forecasting to fault diagnosis. Rather than minimizing per-inverter error, we deliberately train one global Expected Power Model across an entire plant, so that a degraded inverter is measured against an ideal plant-wide baseline rather than against its own compromised history. The resulting deviations, which reduce forecasting accuracy on degraded plants, are precisely the signal that fault diagnosis requires. Building on this principle, our contributions are: (i) a global expected-power modeling approach that structurally avoids anomaly masking; (ii) a physics-informed, rule-based taxonomy that maps deviations to multi-class, action-oriented fault categories; (iii) a Masked Loss Aggregation technique that eliminates Positive Drift Bias when converting detected faults into energy and carbon-emission losses; and (iv) a demonstration that this workflow operates on short-term operational data alone. Using a 34-day public benchmark, we show that practical inverter-level performance monitoring can be established without the multi-month or multi-year records required by many long-term degradation studies, though we note this necessarily precludes the study of long-horizon phenomena such as seasonal effects and soiling cycles.

% =========================
% 2. Related Work
% =========================
Recent studies have increasingly applied Machine Learning to PV forecasting and anomaly detection. Unsupervised Machine Learning (UML) approaches such as One-Class SVM~\cite{scholkopf2001} and Local Outlier Factor~\cite{breunig2000} are widely applied to anomaly detection in PV systems~\cite{mellit2022}. However, as our evaluation in Section~\ref{sec:synthetic_eval} demonstrates, such methods do not incorporate the physical constraints of the DC-to-AC conversion process, and consequently exhibit high false-positive rates when assessed on physically consistent fault data. Other studies emphasize Explainable Artificial Intelligence \cite{arrieta2020} using Extreme Gradient Boosting (XGBoost) \cite{chen2016} and SHapley Additive exPlanations (SHAP) \cite{lundberg2017} for interpretability, but they rarely combine this with a multi-class taxonomy mapped to actionable O\&M thresholds.

Advanced deep learning and machine learning techniques have also been explored for PV fault diagnosis and time-series anomaly detection. For instance, Chen et al.~\cite{chenz2019} proposed a deep residual network operating on I-V curves for array-level fault diagnosis, while Voutsinas et al.~\cite{voutsinas2023} developed ML-based methods for early fault detection. For multivariate time-series anomaly detection broadly, transformer-based architectures such as TranAD~\cite{tuli2022} have achieved state-of-the-art detection performance. On the exact benchmark dataset evaluated in this work, Siddiqa et al.~\cite{siddiqa2025} recently introduced a two-stage transformer framework (SolarTrans) to forecast power generation. However, while these advanced architectures prioritize raw predictive or anomaly detection accuracy, they generally do not provide an action-oriented multi-class taxonomy mapped to specific O\&M dispatch procedures, nor do they quantify energy and carbon loss—critical operational capabilities provided by our proposed framework.

% =========================
% 3. Methodology
% =========================
\section{Methodology}

\subsection{Data and Preprocessing}
Zero-power intervals pose a choice in preprocessing: discarding them avoids division-by-zero when computing conversion efficiency, but in fault detection a zero-power reading under daylight is the most informative signal of total failure. We therefore retain these records and encode undefined efficiency as 0.0.

\subsection{Expected Power Model}
We train an XGBoost \cite{chen2016} regression model to predict the expected AC power output. Crucially, we proactively exclude DC power from the feature set. Including DC power causes target leakage, where the model would merely learn the conversion efficiency rather than the actual impact of environmental variables such as irradiance and temperature. Relying solely on meteorological features ensures the model predicts the ideal power generation, allowing upstream faults to reflect as actionable power losses. Our single global model establishes a plant wide performance baseline, avoiding the anomaly masking inherent to inverter specific local models.

\subsection{Rule-Based Taxonomy}
We implement a multi-class taxonomy designed around actionable O\&M thresholds. The taxonomy identifies categories such as total loss, requiring immediate crew dispatch, and partial loss, requiring inspection within 24 hours. We iteratively refined the taxonomy by removing physically flawed rules; for instance, inverter clipping was removed as its traditional definition contradicted the physical reality of inverter saturation. Table~\ref{tab:taxonomy_thresholds} details the specific detection conditions and corresponding O\&M actions for each class.

\begin{table}[H]
\centering
\caption{Rule-Based Fault Taxonomy: Detection Thresholds and O\&M Actions}
\label{tab:taxonomy_thresholds}
\begin{tabular}{lll}
\toprule
\textbf{Class} & \textbf{Detection Condition} & \textbf{O\&M Action} \\
\midrule
TOTAL\_LOSS & $P_{AC}=0$ and $G>0.2$ & Immediate dispatch \\
PARTIAL\_LOSS & $0 < P_{AC} < 0.5\,\bar{P}_{\text{peer}}$ and $G>0.2$ & Inspect within 24h \\
NORMAL & otherwise & None \\
\bottomrule
\end{tabular}
\end{table}

\noindent where $P_{AC}$ is the measured AC power, $G$ is irradiance (kW/m\textsuperscript{2}), and $\bar{P}_{\text{peer}}$ is the mean AC power of co-located inverters at the same timestamp.

Two design points warrant explanation. First, the PARTIAL\_LOSS rule compares an inverter against the contemporaneous mean of its co-located peers ($\bar{P}_{\text{peer}}$), a relative criterion that complements the absolute plant-wide baseline provided by the global Expected Power Model; the two operate at different granularities and serve as independent corroborating signals rather than a single global comparison. Second, we acknowledge that the TOTAL\_LOSS condition ($P_{AC}=0$ under non-trivial irradiance) closely parallels the mechanism by which total-loss faults are synthetically injected in our evaluation. This overlap is the source of the perfect precision reported in Section~\ref{sec:synthetic_eval} and is discussed as a limitation in Section~\ref{sec:limitations}; the thresholds themselves are grounded in physical operating logic (an inverter producing zero power under sunlight is unambiguously faulted) rather than tuned to the injection procedure.

\subsection{Masked Loss Aggregation}
Quantifying energy loss typically involves a Rectified Linear Unit function taking the maximum of the difference between expected and actual power and zero. However, applying this continuously over time accumulates normal predictive variance into massive phantom losses, a phenomenon we term \textit{positive drift bias}. Formally, given expected power $\hat{y}$ and actual power $y$, continuous accumulation over time $t$ yields $\sum_{t} \max(\hat{y}_t - y_t, 0)$. Because normal stochastic overpredictions are integrated while underpredictions are bounded at zero, the total estimated loss diverges systematically from reality. We introduce Masked Loss Aggregation, which computes power loss exclusively when the taxonomy confirms an anomaly. During normal operation periods, the loss is masked to zero, ensuring high fidelity carbon and energy loss estimations.

\subsection{Explainability}
To build trust with O\&M operators, we integrate SHAP \cite{lundberg2017} to interpret the XGBoost predictions. This allows operators to visualize exactly how ambient temperature and solar irradiation contribute to the expected power target for any given 15 minute interval.

% =========================
% 4. Experimental Setup
% =========================
\section{Experimental Setup}

\subsection{Dataset}
The framework is evaluated on a real world, 34 day dataset collected from two utility scale solar plants in India, comprising 136,476 rows (15 minute intervals) and 33 features. The short timeframe demonstrates performance monitoring on short-horizon operational data. Figure \ref{fig:data_quality} illustrates the data quality before and after preprocessing.

\begin{figure}[H]
    \centering
    \includegraphics[width=0.8\linewidth]{fig_01_data_quality_before_after.png}
    \caption{Data quality and DC correction evidence before and after preprocessing.}
    \label{fig:data_quality}
\end{figure}

\subsection{Physics Informed Synthetic Fault Injection}
To objectively evaluate our rule based taxonomy against UML baselines, we designed a physics informed synthetic fault injection pipeline. When simulating a fault, we ensure data integrity by dynamically recalculating all derived features, such as efficiency and peer average gaps. We also ensure physics integrity. A drop in AC power must be accompanied by a proportional drop in DC power to maintain physical plausibility. Injecting faults without respecting this physical manifold provides an unfair advantage to UML algorithms, which would trivially flag impossible data points.

% =========================
% 5. Results
% =========================
\section{Results}

\subsection{Expected Power Modeling}

\begin{table}[H]
\caption{XGBoost Regression Results}
\begin{center}
\begin{tabular}{lrrrr}
\toprule
\textbf{Plant Split} & \textbf{R\textsuperscript{2}} & \textbf{MAE (kW)} & \textbf{RMSE (kW)} & \textbf{MAPE (\%)} \\
\midrule
Plant 1 Train & 0.9775 & 29.10 & 58.00 & 6.03 \\
Plant 1 Test & 0.9619 & 31.08 & 68.59 & 9.24 \\
Plant 2 Train & 0.5403 & 156.72 & 279.36 & 22.18 \\
Plant 2 Test & 0.5386 & 101.45 & 199.23 & 33.38 \\
\bottomrule
\end{tabular}
\label{tab:xgb_results}
\end{center}
\end{table}

As shown in Table \ref{tab:xgb_results}, the global XGBoost model achieves high accuracy (R\textsuperscript{2} = 0.9619) on the well-maintained Plant 1. For Plant 2, the substantially lower R\textsuperscript{2} (0.5386) is not a symptom of global-model misspecification but a direct measurement of the degradation our baseline is designed to expose. This interpretation is supported by four independent observations. First, the accuracy gap is plant-specific: on the well-maintained Plant 1, our global model attains accuracy comparable to the per-inverter local models of Nahar et al.~\cite{nahar2025} (0.9619 vs.\ 0.98), whereas a large divergence appears only on Plant 2 (0.5386 vs.\ 0.91). A systematic bias inherent to the global architecture would degrade both plants comparably; its confinement to Plant 2 instead points to a genuine site-level effect. Second, this site-level disparity is independently corroborated by Siddiqa et al.~\cite{siddiqa2025}, who applied a two-stage transformer framework (SolarTrans) to the exact same 34-day benchmark dataset and similarly observed a marked performance difference between Plant 1 ($R^2 = 0.9692$) and Plant 2 ($R^2 = 0.7956$). Third, Nahar et al.~\cite{nahar2025}, using localized forecasting on the same plants, independently reported that Plant 2 exhibits a bimodal inverter-performance distribution with a subset of inverters consistently underperforming (R\textsuperscript{2} below 80\%), which they attribute to localized shading, persistent soiling, or hardware degradation. The chronic underperformance our global baseline surfaces is thus consistent with physically documented site conditions. Fourth, the near-identical Train and Test R\textsuperscript{2} (0.5403 vs.\ 0.5386) rules out overfitting as the source of this variance. To visually demonstrate this, Figure \ref{fig:residual_plant2} plots the time-series residuals (Actual - Expected AC Power) for Plant 2, highlighting how the identified anomaly clusters correlate directly with systematic power deficits. Figure \ref{fig:regression} visualizes the model's actual versus predicted performance.

\begin{figure}[H]
    \centering
    \includegraphics[width=0.8\linewidth]{fig_residual_plant2.png}
    \caption{Plant 2 time-series residuals highlighting systematic degradation.}
    \label{fig:residual_plant2}
\end{figure}

\begin{figure}[H]
    \centering
    \includegraphics[width=0.8\linewidth]{fig_03_regression_actual_vs_predicted.png}
    \caption{XGBoost global model regression actual versus predicted.}
    \label{fig:regression}
\end{figure}



Feature importance is quantified using SHAP values, demonstrating the dominance of irradiance over temperature in the expected power generation, as shown in Figure \ref{fig:shap_feature}.

\begin{figure}[H]
    \centering
    \includegraphics[width=0.8\linewidth]{fig_04_shap_feature_importance.png}
    \caption{SHAP feature importance for the Expected Power Model.}
    \label{fig:shap_feature}
\end{figure}

\subsection{Anomaly Distribution}
Table \ref{tab:anomaly_dist} details the frequency of detected faults. Thirty-six intervals exhibit an AC to DC ratio marginally exceeding unity, a known artefact of 15-minute interval averaging. The vast majority of instances are normal operations. It should be explicitly noted that classes such as thermal degrade and sensor drift are purely theoretical within the scope of this study; they were designed for long-term monitoring scenarios, and zero empirical instances were observed in the 34-day evaluation window. Consequently, they are omitted from further empirical validation. Note that Table~\ref{tab:anomaly_dist} reports counts evaluated over the entire dataset (136,476 total intervals, including night-time periods where output is zero by definition and classified as \texttt{NORMAL}), whereas Figure~\ref{fig:anomaly_dist} isolates daytime operational hours (72,904 intervals). The resulting difference in \texttt{NORMAL} counts (132,282 full-set vs.\ 72,904 daytime) reflects this filtering convention rather than a numerical discrepancy.

% Figure \ref{fig:anomaly_dist} provides a visual breakdown of this taxonomy distribution.

\begin{table}[H]
\caption{Anomaly Taxonomy Distribution}
\begin{center}
\begin{tabular}{lrr}
\toprule
\textbf{Class} & \textbf{Count} & \textbf{\%} \\
\midrule
\texttt{NORMAL} & 132,282 & 96.93\% \\
\texttt{TOTAL\_LOSS} & 3,821 & 2.80\% \\
\texttt{PARTIAL\_LOSS} & 373 & 0.27\% \\
\bottomrule
\end{tabular}
\label{tab:anomaly_dist}
\end{center}
\end{table}

% \begin{figure}[H]
%     \centering
%     \includegraphics[width=0.8\linewidth]{fig_06_anomaly_class_distribution.png}
%     \caption{Distribution of multi class anomaly taxonomy.}
%     \label{fig:anomaly_dist}
% \end{figure}

\subsection{Synthetic Fault Evaluation}
\label{sec:synthetic_eval}
To avoid circularity between hyperparameter selection and performance
reporting, we adopted a two-stage evaluation protocol. In the first
stage, we performed a grid search on an independent validation set
consisting of ten synthetically injected seeds (Seeds 1--10), never
used in final performance estimation. The search space was:
Isolation Forest $n\_estimators \in \{50, 100, 200, 300, 500\}$ with
$\mathrm{max\_samples}$ fixed at the sklearn default (\texttt{auto}
$= \min(256, n)$); One-Class SVM $\nu \in \{0.01, 0.02, 0.03, 0.05,
0.1\}$; and LOF $k \in \{10, 15, 20, 30, 50\}$. For each candidate,
we fit the model on the full injected dataset and evaluated the
maximum F1-score along the Precision--Recall curve (i.e., the optimal
decision threshold). Hyperparameter selection was based on the mean
F1-score across all ten validation seeds.

Note that for Isolation Forest, \texttt{contamination} was not included
in the search space. The \texttt{contamination} parameter only shifts
the decision threshold in \texttt{predict()}, and has no effect on the
anomaly scores returned by \texttt{score\_samples()}, which we use
directly for PR curve evaluation. The genuinely influential parameters
are \texttt{n\_estimators} (ensemble diversity) and
\texttt{max\_samples} (sub-sample size). Based on Liu et al.\
\cite{liu2012}, sub-sample sizes beyond 256 provide negligible
improvement and can increase swamping and masking artefacts; accordingly
we fix \texttt{max\_samples = auto}. The validation F1-score was
observed to approach saturation at $n\_estimators = 500$ (Avg F1 =
0.289 vs.\ 0.284 at $n\_estimators = 300$, $\Delta = 0.005$), so 500
was selected as the optimal value.

In the second stage, the three optimal hyperparameter sets---Isolation
Forest ($n\_estimators = 500$, $\mathrm{max\_samples} = \texttt{auto}$),
One-Class SVM ($\nu = 0.02$), and LOF ($k = 10$)---were applied to a
completely independent test set of 20 seeds (Seeds 11--30). For each
model and seed, the optimal decision threshold maximising F1 on the
Precision--Recall curve was identified independently, and Cohen's
$\kappa$ was computed at this same threshold. Performance is reported
as mean $\pm$ standard deviation across all 20 seeds (Table~\ref{tab:synthetic_faults}).

\begin{table}[H]
\caption{Synthetic Fault Evaluation (Physics-Informed, N=20 Seeds). Precision, Recall, F1, and Cohen's $\kappa$ are computed at the optimal Precision--Recall threshold per seed; mean $\pm$ std reported over Seeds 11--30.}
\begin{center}
\begin{tabular}{lrrrr}
\toprule
\textbf{Method} & \textbf{Precision} & \textbf{Recall} & \textbf{F1} & \textbf{Kappa} \\
\midrule
Rule-Based Taxonomy & $1.00 \pm 0.00$ & $0.58 \pm 0.06$ & $0.74 \pm 0.05$ & $0.74 \pm 0.05$ \\
Isolation Forest    & $0.32 \pm 0.04$ & $0.29 \pm 0.02$ & $0.30 \pm 0.02$ & $0.30 \pm 0.02$ \\
One-Class SVM       & $0.17 \pm 0.04$ & $0.52 \pm 0.11$ & $0.25 \pm 0.03$ & $0.24 \pm 0.03$ \\
LOF                 & $0.15 \pm 0.06$ & $0.39 \pm 0.10$ & $0.21 \pm 0.05$ & $0.21 \pm 0.05$ \\
\bottomrule
\end{tabular}
\label{tab:synthetic_faults}
\end{center}
\end{table}

Table~\ref{tab:synthetic_faults} compares detection performance across all methods. The rule-based taxonomy achieves a perfect Precision of $1.00 \pm 0.00$, ensuring zero false positive alarms. The Recall of $0.58 \pm 0.06$ reflects the intentional actionable threshold design. Class-specific recall further confirmed high sensitivity for catastrophic events, achieving $0.69 \pm 0.07$ for TOTAL\_LOSS and $0.66 \pm 0.04$ for PARTIAL\_LOSS (N=20, Seeds 11--30), demonstrating its utility for prioritizing the most operationally severe anomalies. Across all evaluated methods, Cohen's $\kappa$ closely tracks the F1 score due to the severe class imbalance in the synthetic evaluation dataset, which reduces chance agreement terms and causes $\kappa$ to converge numerically toward F1. Even after rigorous hyperparameter tuning on an independent validation set, all UML baselines demonstrate substantially lower performance, confirming the superiority of the physics-informed rule-based approach on physically realistic fault data. The Isolation Forest achieves a stable Precision of $0.32 \pm 0.04$ (notably reduced variance compared to untuned runs), while LOF and One-Class SVM exhibit higher Recall at the cost of lower Precision. Figure~\ref{fig:pr_curve} presents the Precision--Recall curves, and Figure~\ref{fig:baseline_comparison} illustrates this performance gap graphically.

\begin{figure}[H]
    \centering
    \includegraphics[width=0.8\linewidth]{fig_06_pr_curves.png}
    \caption{Precision-Recall curves for anomaly detection baselines. UML curves (IF, OCSVM, LOF) are illustrative single-seed curves from Seed 11; the Rule-Based operating point represents the aggregate mean over N=20 test seeds (Seeds 11--30).}
    \label{fig:pr_curve}
\end{figure}



\begin{figure}[H]
    \centering
    \includegraphics[width=0.8\linewidth]{fig_08_baseline_comparison.png}
    \caption{Baseline comparison between Rule-Based Taxonomy and UML models (Mean $\pm$ Std, N=20 seeds, Seeds 11--30). IF: $n\_estimators=500$, max\_samples=auto; OCSVM: $\nu=0.02$; LOF: $k=10$.}
    \label{fig:baseline_comparison}
\end{figure}

\subsection{Carbon and Energy Loss Quantification}

\begin{table}[H]
\caption{Carbon \& Energy Loss (Bootstrap 95\% CI)}
\begin{center}
\begin{tabular}{lrrr}
\toprule
\textbf{Plant} & \textbf{Energy Loss (kWh)} & \textbf{CO2 Loss (kg)} & \textbf{Loss Rate (\%)} \\
\midrule
Plant 1 & 16,452.13 & 11,565.85 & 0.31\% \\
Plant 2 & 622,853.97 & 437,866.34 & 15.25\% \\
\textbf{TOTAL} & \textbf{639,306.10} [CI: 525k -- 748k] & \textbf{449,432.19} [CI: 369k -- 526k] & -- \\
\bottomrule
\end{tabular}
\label{tab:carbon_loss}
\end{center}
\end{table}

Table \ref{tab:carbon_loss} highlights the plant level energy and carbon losses. Leveraging Masked Loss Aggregation, we isolated exactly 639,306.10 kWh of recoverable energy loss, translating to 449,432.19 kg of avoided CO2. The carbon footprint estimation is based on the average Indian Grid Emission Factor (0.703 kgCO2/kWh) published by the Central Electricity Authority (CEA) for FY2020-21~\cite{cea2021}. The energy loss breakdown across different anomaly classes is detailed in Figure \ref{fig:energy_loss}. To illustrate the practical value of the system, a specific case study timeline is shown in Figure \ref{fig:case_timeline}.



\begin{figure}[H]
    \centering
    \includegraphics[width=0.8\linewidth]{fig_14_energy_loss_by_class.png}
    \caption{Recoverable energy loss breakdown by anomaly class.}
    \label{fig:energy_loss}
\end{figure}

\begin{figure}[H]
    \centering
    \includegraphics[width=0.8\linewidth]{fig_11_case1_timeline.png}
    \caption{Case study timeline showing detected actionable anomalies over time.}
    \label{fig:case_timeline}
\end{figure}

% =========================
% 6. Discussion
% =========================
\section{Discussion}

The results validate several key insights that underpin this framework. Prioritizing a global model over local models prevents anomaly masking, ensuring that degraded inverters are evaluated against an ideal plant wide baseline. This design choice is fundamental for identifying chronic underperformance rather than simply modeling it as the new normal.

The introduction of Masked Loss Aggregation effectively neutralizes positive drift bias. Without this mechanism, the natural variance of the XGBoost model would continuously accumulate phantom losses, destroying the credibility of the O\&M reports. By only unmasking losses when the taxonomy detects an anomaly, the calculations remain tightly coupled to actual, recoverable events.

Physical consistency in synthetic evaluation materially affects the conclusions drawn. When faults are injected respecting the physical limits of the DC-to-AC conversion process, and all derived features are recalculated accordingly, the unsupervised baselines produce substantially more false positives than under naive injection. Injecting faults that violate DC-AC consistency would allow these methods to flag physically impossible points trivially, inflating their apparent performance. Removing rules that are not physically grounded likewise ensures that O\&M teams are dispatched only for genuine, recoverable losses.

\subsection{Limitations}
\label{sec:limitations}
Two limitations of our evaluation must be stated explicitly. First, the synthetic fault evaluation exhibits partial circularity: the total-loss detection rule (AC power near zero under non-trivial irradiance) closely mirrors the mechanism by which total-loss faults are injected. Consequently, the perfect precision reported in Table~\ref{tab:synthetic_faults} should be interpreted as a consistency check confirming that the rule set correctly recovers faults conforming to its own physical assumptions, rather than as evidence of detection capability on independent, real-world faults. Second, no field-verified fault records (e.g., O\&M logs or repair reports) were available to validate detections against ground truth; all reported detection metrics are computed on synthetically injected faults. Validation on real annotated faults remains important future work.

The energy and carbon loss estimates (639 MWh, 449 tCO\textsubscript{2}) depend on the global expected-power baseline, which does not distinguish genuine underperformance from legitimate inter-inverter variability (e.g., differences in siting, capacity, or connection). A portion of the estimated Plant 2 loss (15.25\%) may therefore reflect systematic model bias rather than recoverable loss. While Masked Loss Aggregation removes the random accumulation of positive prediction noise, it does not correct for systematic baseline bias. These figures should thus be read as an upper-bound estimate of recoverable loss, pending validation against actual recovered-energy or repair records.

% =========================
% 7. Conclusion
% =========================
\section{Conclusion}

This paper presented an explainable, data-efficient framework tailored for solar PV monitoring. By combining a global XGBoost model with a physics-informed rule taxonomy and Masked Loss Aggregation, the framework successfully mitigates anomaly masking, positive drift bias, and alert fatigue. Validated on a short 34-day public benchmark, the system identified over 639 MWh of recoverable losses, representing nearly 450 metric tons of avoided CO\textsubscript{2}, while providing an upper-bound estimate of recoverable energy loss. Future work will explore dynamic thresholding and adaptive taxonomies for long-term monitoring across multi-year operational datasets.

% =========================
\clearpage

% References
% =========================
\begingroup
\setlength{\parskip}{0pt}
\begin{thebibliography}{99}
\setlength{\itemsep}{0pt}

\bibitem{nahar2025}
A.~Nahar, R.~A.~M.~Rudro, M.~F.~A.~Al~Sohan, M.~H.~Uddin, and L.~Kumar. Forecasting solar photovoltaic power generation: A machine learning time series model approach. \emph{International Journal of Energy Research}, 2025:4092367, 2025.
\href{https://doi.org/10.1155/er/4092367}{doi:10.1155/er/4092367}.

\bibitem{siddiqa2025}
A.~Siddiqa, N.~Rana, W.~Z.~Khan, F.~Jeribi, and A.~Tahir. Explaining solar forecasts with generative AI: A two-stage framework combining transformers and LLMs. \emph{PLOS ONE}, 20(9):e0331516, 2025.
\href{https://doi.org/10.1371/journal.pone.0331516}{doi:10.1371/journal.pone.0331516}.

\bibitem{mellit2022}
A.~Mellit and S.~Kalogirou. Assessment of machine learning and ensemble methods for fault diagnosis of photovoltaic systems. \emph{Renewable Energy}, 184:1074--1090, 2022.
\href{https://doi.org/10.1016/j.renene.2021.11.125}{doi:10.1016/j.renene.2021.11.125}.

\bibitem{voutsinas2023}
S.~Voutsinas, D.~Karolidis, I.~Voyiatzis, et~al. Development of a machine-learning-based method for early fault detection in photovoltaic systems. \emph{Journal of Engineering and Applied Science}, 70:27, 2023.
\href{https://doi.org/10.1186/s44147-023-00200-0}{doi:10.1186/s44147-023-00200-0}.

\bibitem{chenz2019}
Z.~Chen, Y.~Chen, L.~Wu, S.~Cheng, and P.~Lin. Deep residual network based fault detection and diagnosis of photovoltaic arrays using current-voltage curves and ambient conditions. \emph{Energy Conversion and Management}, 198:111793, 2019.
\href{https://doi.org/10.1016/j.enconman.2019.111793}{doi:10.1016/j.enconman.2019.111793}.

\bibitem{chen2016}
T.~Chen and C.~Guestrin. XGBoost: A scalable tree boosting system. In \emph{Proceedings of the 22nd ACM SIGKDD International Conference on Knowledge Discovery and Data Mining}, pages 785--794, 2016.
\href{https://doi.org/10.1145/2939672.2939785}{doi:10.1145/2939672.2939785}.

\bibitem{lundberg2017}
S.~M.~Lundberg and S.-I.~Lee. A unified approach to interpreting model predictions. In \emph{Advances in Neural Information Processing Systems 30}, pages 4765--4774, 2017.
\href{https://doi.org/10.48550/arXiv.1705.07874}{doi:10.48550/arXiv.1705.07874}.

\bibitem{arrieta2020}
A.~B.~Arrieta, N.~D{\'i}az-Rodr{\'i}guez, J.~Del~Ser, A.~Bennetot, S.~Tabik, A.~Barbado, S.~Garc{\'i}a, S.~Gil-L{\'o}pez, D.~Molina, R.~Benjamins, R.~Chatila, and F.~Herrera. Explainable Artificial Intelligence (XAI): Concepts, taxonomies, opportunities and challenges toward responsible AI. \emph{Information Fusion}, 58:82--115, 2020.
\href{https://doi.org/10.1016/j.inffus.2019.12.012}{doi:10.1016/j.inffus.2019.12.012}.

\bibitem{tuli2022}
S.~Tuli, G.~Casale, and N.~R.~Jennings. TranAD: Deep transformer networks for anomaly detection in multivariate time series data. \emph{Proceedings of the VLDB Endowment}, 15(6):1201--1214, 2022.
\href{https://doi.org/10.14778/3514061.3514067}{doi:10.14778/3514061.3514067}.

\bibitem{liu2012}
F.~T.~Liu, K.~M.~Ting, and Z.-H.~Zhou. Isolation-based anomaly detection. \emph{ACM Transactions on Knowledge Discovery from Data}, 6(1):3, 2012.
\href{https://doi.org/10.1145/2133360.2133363}{doi:10.1145/2133360.2133363}.

\bibitem{breunig2000}
M.~M.~Breunig, H.-P.~Kriegel, R.~T.~Ng, and J.~Sander. LOF: Identifying density-based local outliers. \emph{ACM SIGMOD Record}, 29(2):93--104, 2000.
\href{https://doi.org/10.1145/342009.335388}{doi:10.1145/342009.335388}.

\bibitem{scholkopf2001}
B.~Sch{\"o}lkopf, J.~C.~Platt, J.~Shawe-Taylor, A.~J.~Smola, and R.~C.~Williamson. Estimating the support of a high-dimensional distribution. \emph{Neural Computation}, 13(7):1443--1471, 2001.
\href{https://doi.org/10.1162/089976601750264965}{doi:10.1162/089976601750264965}.

\bibitem{cea2021}
Central Electricity Authority (CEA). \emph{CO2 Baseline Database for the Indian Power Sector, User Guide Version 17.0}. Ministry of Power, Government of India, New Delhi, 2021.
\href{https://cea.nic.in/cdm-co2-baseline-database/}{https://cea.nic.in/cdm-co2-baseline-database/}.

\end{thebibliography}
\endgroup

\end{document}
